% GENERATED FILE. Edit canonical lessons, then run npm run generate:books.
% canonical-source-hash: fnv1a64:a74ddaf59ddb8ccd
% canonical-lessons: PA-C109-mahar, PA-C109-jankari, PA-C109-sakna, PA-C109-kol, PA-C109-takat

\chapter{Expert, Knowledge, To be able to, Near, Strength}
\label{ch:pa-expert-knowledge-to-be-able-to-near-strength}
\hlchaptermodality{109}

\begin{chapteropening}
\textbf{By the end of this chapter:} \emph{I can say an expert, knowledge, to be able to, near and strength.}

Complete the last lesson of chapter 109: I can say an expert, knowledge, to be able to, near and strength.
\end{chapteropening}

\section[\texorpdfstring{\pa{ਮਾਹਰ} (māhar)}{ਮਾਹਰ (māhar)}]{\pa{ਮਾਹਰ} (māhar) — an expert}

\label{lesson:PA-C109-mahar}



\begin{quote}
Before the new one: say the Punjabi for like that, then the Punjabi for other.
\end{quote}

\subsection*{You'll want to know: \pa{ਮਾਹਰ}}
\textbf{\pa{ਮਾਹਰ}} — \emph{māhar} — \textquotedblleft{}an expert\textquotedblright{}.

\begin{grammarlens}[title={What you've built}]
One more word for this chapter.
\end{grammarlens}

\subsection*{Your turn}
\begin{itemize}
\raggedright
  \item \emph{Say it:} \emph{māhar}
  \item \emph{Say it:} \emph{māhar}, once more
  \item \emph{Say it:} say \emph{māhar}
  \item \emph{Recall:} say the Punjabi for like that, then the Punjabi for other, then say \emph{māhar} again
\end{itemize}

\begin{tcolorbox}[breakable,colback=teal!4,colframe=teal!35!black,title={Before you move on}]
What does \pa{ਮਾਹਰ} mean? (\textquotedblleft{}An expert\textquotedblright{}.) Say it once more.
\end{tcolorbox}
\section[\texorpdfstring{\pa{ਜਾਣਕਾਰੀ} (jāṇkārī)}{ਜਾਣਕਾਰੀ (jāṇkārī)}]{\pa{ਜਾਣਕਾਰੀ} (jāṇkārī) — knowledge, information}

\label{lesson:PA-C109-jankari}



\begin{quote}
Before the new one: say the Punjabi for other, then the Punjabi for an expert.
\end{quote}

\subsection*{You'll want to know: \pa{ਜਾਣਕਾਰੀ}}
\textbf{\pa{ਜਾਣਕਾਰੀ}} — \emph{jāṇkārī} — \textquotedblleft{}knowledge, information\textquotedblright{}.

\begin{grammarlens}[title={What you've built}]
One more word for this chapter.
\end{grammarlens}

\subsection*{Your turn}
\begin{itemize}
\raggedright
  \item \emph{Say it:} \emph{jāṇkārī}
  \item \emph{Say it:} \emph{jāṇkārī}, once more
  \item \emph{Say it:} say \emph{jāṇkārī}
  \item \emph{Recall:} say the Punjabi for other, then the Punjabi for an expert, then say \emph{jāṇkārī} again
\end{itemize}

\begin{tcolorbox}[breakable,colback=teal!4,colframe=teal!35!black,title={Before you move on}]
What does \pa{ਜਾਣਕਾਰੀ} mean? (\textquotedblleft{}Knowledge, information\textquotedblright{}.) Say it once more.
\end{tcolorbox}
\section[\texorpdfstring{\pa{ਸਕਣਾ} (sakṇā)}{ਸਕਣਾ (sakṇā)}]{\pa{ਸਕਣਾ} (sakṇā) — to be able to, can}

\label{lesson:PA-C109-sakna}



\begin{quote}
Before the new one: say the Punjabi for an expert, then the Punjabi for knowledge.
\end{quote}

\subsection*{You'll want to know: \pa{ਸਕਣਾ}}
\textbf{\pa{ਸਕਣਾ}} — \emph{sakṇā} — \textquotedblleft{}to be able to, can\textquotedblright{}.

\begin{grammarlens}[title={What you've built}]
One more word for this chapter.
\end{grammarlens}

\subsection*{Your turn}
\begin{itemize}
\raggedright
  \item \emph{Say it:} \emph{sakṇā}
  \item \emph{Say it:} \emph{sakṇā}, once more
  \item \emph{Say it:} say \emph{sakṇā}
  \item \emph{Recall:} say the Punjabi for an expert, then the Punjabi for knowledge, then say \emph{sakṇā} again
\end{itemize}

\begin{tcolorbox}[breakable,colback=teal!4,colframe=teal!35!black,title={Before you move on}]
What does \pa{ਸਕਣਾ} mean? (\textquotedblleft{}To be able to, can\textquotedblright{}.) Say it once more.
\end{tcolorbox}
\section[\texorpdfstring{\pa{ਕੋਲ} (kol)}{ਕੋਲ (kol)}]{\pa{ਕੋਲ} (kol) — near, in one's possession}

\label{lesson:PA-C109-kol}



\begin{quote}
Before the new one: say the Punjabi for knowledge, then the Punjabi for to be able to.
\end{quote}

\subsection*{You'll want to know: \pa{ਕੋਲ}}
\textbf{\pa{ਕੋਲ}} — \emph{kol} — \textquotedblleft{}near, in one's possession\textquotedblright{}.

\begin{grammarlens}[title={What you've built}]
One more word for this chapter.
\end{grammarlens}

\subsection*{Your turn}
\begin{itemize}
\raggedright
  \item \emph{Say it:} \emph{kol}
  \item \emph{Say it:} \emph{kol}, once more
  \item \emph{Say it:} say \emph{kol}
  \item \emph{Recall:} say the Punjabi for knowledge, then the Punjabi for to be able to, then say \emph{kol} again
\end{itemize}

\begin{tcolorbox}[breakable,colback=teal!4,colframe=teal!35!black,title={Before you move on}]
What does \pa{ਕੋਲ} mean? (\textquotedblleft{}Near, in one's possession\textquotedblright{}.) Say it once more.
\end{tcolorbox}
\section[\texorpdfstring{\pa{ਤਾਕਤ} (tākat)}{ਤਾਕਤ (tākat)}]{\pa{ਤਾਕਤ} (tākat) — strength}

\label{lesson:PA-C109-takat}



\begin{quote}
Before the new one: say the Punjabi for to be able to, then the Punjabi for near.
\end{quote}

\subsection*{You'll want to know: \pa{ਤਾਕਤ}}
\textbf{\pa{ਤਾਕਤ}} — \emph{tākat} — \textquotedblleft{}strength\textquotedblright{}.

\begin{grammarlens}[title={What you've built}]
One more word for this chapter.
\end{grammarlens}

\subsection*{Your turn}
\begin{itemize}
\raggedright
  \item \emph{Say it:} \emph{tākat}
  \item \emph{Say it:} \emph{tākat}, once more
  \item \emph{Say it:} say \emph{tākat}
  \item \emph{Recall:} say the Punjabi for to be able to, then the Punjabi for near, then say \emph{tākat} again
\end{itemize}

\begin{tcolorbox}[breakable,colback=teal!4,colframe=teal!35!black,title={Before you move on}]
What does \pa{ਤਾਕਤ} mean? (\textquotedblleft{}Strength\textquotedblright{}.) Say it once more.
\end{tcolorbox}
